% !TeX program = xelatex
% ==========================================================================
% Звіт з лабораторної роботи — ДВВ «Основи прикладної фізики» (1.26-27), УКУ.
% Викладач: Володимир Коломієць.
% ==========================================================================
\documentclass[ukrainian]{ucaphyslab}

% --------------------------------------------------------------------------
% Метадані титульної сторінки
% --------------------------------------------------------------------------
\labtitle{Маятник Обербека}
\labsubtitle{Лабораторна робота № 2}

\student{Довбенчук Олена}
\student{Кульбаба Андрій}
\student{Магарита Ія}
\student{Пастарнак Марко}

\teacher{Володимир Коломієць}

\makeatletter
\renewcommand{\ucu@labelstudent}{Студенти}
\renewcommand{\ucu@labelteacher}{Викладач}
\makeatother

\labdate{2026-10-14}

\usepackage{float}
\usepackage{pgfplots}
\pgfplotsset{compat=1.18}
\sisetup{uncertainty-mode=separate}
\definecolor{series1}{HTML}{2A78D6}
\definecolor{series2}{HTML}{EB6834}

\begin{document}
\maketitlepage

% ==========================================================================
\section{Мета роботи}
% ==========================================================================
\begin{enumerate}
  \item Дослідити динаміку обертального руху абсолютно твердого тіла на прикладі хрестоподібного маятника Обербека.
  \item Експериментально перевірити основне рівняння динаміки обертального руху (рівняння моментів) навколо нерухомої осі двома методами: за залежністю кутового прискорення $\beta$ від моменту зовнішньої сили $mgr$ при сталому моменті інерції ($I = \text{const}$) та за залежністю відношення $g/a$ від безрозмірного параметра $(R/r)^2$ при змінному моменті інерції.
  \item Визначити момент сили тертя в підшипниках осі $M_T$, сумарний момент інерції навантаженого маятника $I$ та власний момент інерції ненавантаженого маятника $I_0$.
\end{enumerate}

% ==========================================================================
\section{Прилади та обладнання}
% ==========================================================================
\begin{itemize}
  \item \textbf{Маятник Обербека} — втулка з чотирма взаємно перпендикулярними спицями, двома шківами радіусів $r_1 = \SI{20}{\milli\meter}$ і $r_2 = \SI{30}{\milli\meter}$, чотирма рухомими муфточками та платформою для тягарців на нитці.
  \item \textbf{Вертикальна мірна лінійка} — ціна поділки $\SI{1}{\milli\meter}$, інструментальна похибка $\Delta h = \SI{0.5}{\milli\meter}$ (з урахуванням точності фіксації початкового положення платформи прийнято $\Delta h = \SI{1}{\milli\meter}$).
  \item \textbf{Штангенциркуль} — ціна поділки ноніуса $\SI{0.05}{\milli\meter}$.
  \item \textbf{Електронні лабораторні терези} — ціна поділки (дискретність) $\SI{0.01}{\gram}$, похибка зважування $\Delta m = \SI{0.01}{\gram}$.
  \item \textbf{Секундоміри} (вимірювання проводилися паралельно кількома спостерігачами для мінімізації суб'єктивної похибки реакції) — ціна поділки $\SI{0.01}{\second}$.
\end{itemize}

% ==========================================================================
\section{Теоретичні відомості}
% ==========================================================================
Маятник Обербека складається з чотирьох спиць, закріплених на втулці під прямим кутом, уздовж яких можна переміщувати чотири циліндричні муфточки на однакову відстань $R$ від осі обертання. На горизонтальній осі втулки жорстко закріплені два шківи радіусів $r_1$ та $r_2$. На один зі шківів намотується нерозтяжна нитка, до вільного кінця якої прикріплена платформа з тягарцями сумарної маси $m$.

Рух системи описується другим законом Ньютона для поступального руху платформи з тягарцями, рівнянням моментів для обертального руху маятника відносно горизонтальної осі та рівнянням кінематичного зв'язку (за умови відсутності проковзування нитки):
\begin{equation}
  \begin{cases}
    I\beta = F_{\text{н}}r - M_T, \\
    ma = mg - F_{\text{н}}, \\
    a = \beta r,
  \end{cases}
  \label{eq:system}
\end{equation}
де $I$ — момент інерції маятника відносно осі обертання; $\beta$ — кутове прискорення маятника; $a$ — лінійне прискорення платформи; $F_{\text{н}}$ — сила натягу нитки; $M_T$ — гальмівний момент сили тертя в підшипниках осі; $g = \SI{9.81}{\meter\per\second\squared}$ — прискорення вільного падіння.

Розв'язуючи систему \eqref{eq:system} відносно лінійного прискорення $a$, отримуємо:
\begin{equation}
  a = \frac{(mgr^2 - M_T r)/I}{1 + mr^2/I}.
  \label{eq:a_exact}
\end{equation}
Оскільки в умовах експерименту доданок $mr^2/I \ll 1$ (приблизно на два порядки менший за одиницю, що відповідає повільному падінню платформи $a \ll g$), ним у знаменнику \eqref{eq:a_exact} можна знехтувати. Тоді основне робоче рівняння для кутового прискорення набуває лінійного вигляду відносно моменту сили тяжіння $mgr$:
\begin{equation}
  \beta = \frac{a}{r} = \frac{1}{I}(mgr) - \frac{M_T}{I}.
  \label{eq:beta_linear}
\end{equation}

Оскільки рух платформи зі стану спокою є рівноприскореним, прискорення $a$ та $\beta$ визначаються через висоту падіння $h$ і час руху $t$:
\begin{equation}
  a = \frac{2h}{t^2}, \qquad \beta = \frac{2h}{r t^2}.
  \label{eq:kinematics}
\end{equation}

За теоремою Штейнера та властивістю адитивності сумарний момент інерції маятника складається з моменту інерції ненавантаженої хрестовини зі шківами $I_0$ та моменту інерції чотирьох муфточок (у нашій установці — дві муфточки масою $m_s$ та дві масою $m_b$, розташовані симетрично на відстані $R$ від осі):
\begin{equation}
  I = I_0 + \sum_{i=1}^{4} m_{\text{м},i} R^2 = I_0 + 2(m_s + m_b)R^2.
  \label{eq:I_total}
\end{equation}

Якщо обрати масу вантажу достатньо великою ($m \gg m_0$, де $m_0$ — мінімальна маса, за якої маятник долає тертя спокою, так що $mgr \gg M_T$), моментом сили тертя $M_T$ у \eqref{eq:beta_linear} можна знехтувати. Підставляючи \eqref{eq:I_total} у \eqref{eq:beta_linear}, одержуємо лінійну залежність безрозмірного відношення $g/a$ від $(R/r)^2$:
\begin{equation}
  \frac{g}{a} = \frac{I_0}{mr^2} + \frac{2(m_s + m_b)}{m}\left(\frac{R}{r}\right)^2.
  \label{eq:ga_linear}
\end{equation}

% ==========================================================================
\section{Хід роботи та результати вимірювань}
% ==========================================================================
Перед початком основних серій дослідів було виміряно геометричні та масові параметри установки (Табл.~\ref{tab:constants}). Початкова маса важкої платформи становила $M = \SI{102.00}{\gram}$. Шляхом поступового довантаження платформи дрібними важками було знайдено порогові значення додаткової маси $m_{0,1}$ та $m_{0,2}$ (і відповідні повні маси $m_{\text{повна},1} = M + m_{0,1}$ та $m_{\text{повна},2} = M + m_{0,2}$), при яких маятник починає самостійно обертатися при намотуванні нитки на шківи $r_1$ та $r_2$.

\begin{table}[H]
  \centering
  \begin{tabular}{@{}llcl@{}}
    \toprule
    Назва величини & Позначення & Значення & Одиниці \\
    \midrule
    Маса початкова (важка платформа) & $M$ & \num{102.00} & \si{\gram} \\
    Додаткова маса подолання тертя (шків $r_1$) & $m_{0,1}$ & \num{38.80} & \si{\gram} \\
    Додаткова маса подолання тертя (шків $r_2$) & $m_{0,2}$ & \num{26.24} & \si{\gram} \\
    Повна порогова маса руху (шків $r_1$) & $m_{\text{повна},1}$ & \num{140.80} & \si{\gram} \\
    Повна порогова маса руху (шків $r_2$) & $m_{\text{повна},2}$ & \num{128.24} & \si{\gram} \\
    Фіксована висота падіння платформи & $h$ & \num{107.6} & \si{\centi\meter} \\
    Радіус меншого шківа & $r_1$ & \num{20.0} & \si{\milli\meter} \\
    Радіус більшого шківа & $r_2$ & \num{30.0} & \si{\milli\meter} \\
    Фіксована відстань від муфт до осі (Дослід 1) & $R_{\text{фікс}}$ & \num{16.21} & \si{\centi\meter} \\
    Фіксована маса вантажу (Дослід 2) & $m_{\text{фікс}}$ & \num{373.0} & \si{\gram} \\
    Маса меншої муфточки (2 шт.) & $m_s$ & \num{100.77} & \si{\gram} \\
    Маса більшої муфточки (2 шт.) & $m_b$ & \num{125.34} & \si{\gram} \\
    \bottomrule
  \end{tabular}
  \caption{Вихідні константи та параметри експериментальної установки}
  \label{tab:constants}
\end{table}

\subsection{Дослід 1. Залежність кутового прискорення від моменту сили ($I = \text{const}$)}
Муфточки зафіксували на однаковій відстані $R_{\text{фікс}} = \SI{16.21}{\centi\meter}$ від осі обертання. Для кожного з двох шківів ($r_1 = \SI{20}{\milli\meter}$ та $r_2 = \SI{30}{\milli\meter}$) послідовно встановлювали 5 різних значень маси падаючого вантажу $m$. Кожен запуск з висоти $h = \SI{1.076}{\meter}$ повторювали по 5 разів (спроби 1--5), причому в кожній спробі час паралельно фіксували три спостерігачі для усунення грубих похибок. Усереднені результати спроб та розраховані за формулами \eqref{eq:kinematics} величини наведено в Табл.~\ref{tab:exp1}.

\begin{table}[H]
  \centering
  \small
  \begin{tabular}{@{}cccccccccc@{}}
    \toprule
    $m$, \si{\gram} & $\bar{t}_1$, \si{\second} & $\bar{t}_2$, \si{\second} & $\bar{t}_3$, \si{\second} & $\bar{t}_4$, \si{\second} & $\bar{t}_5$, \si{\second} & $\bar{t} \pm \Delta t$, \si{\second} & $mgr$, \si{\newton\meter} & $a$, \si{\meter\per\second\squared} & $\beta \pm \Delta\beta$, \si{\radian\per\second\squared} \\
    \midrule
    \multicolumn{10}{c}{\textbf{Менший шків: $r_1 = \SI{20}{\milli\meter}$}} \\
    \midrule
    373.97 & 7.05 & 8.66 & 10.29 & 9.53 & 10.64 & $9.23 \pm 0.64$ & 0.0734 & 0.0253 & $1.26 \pm 0.18$ \\
    411.21 & 8.06 & 11.55 & 8.79 & 7.00 & 6.47 & $8.37 \pm 0.89$ & 0.0807 & 0.0307 & $1.53 \pm 0.33$ \\
    449.30 & 7.18 & 7.45 & 6.81 & 6.21 & 6.28 & $6.79 \pm 0.25$ & 0.0882 & 0.0467 & $2.34 \pm 0.17$ \\
    508.37 & 5.91 & 5.57 & 6.73 & 5.80 & 6.47 & $6.09 \pm 0.22$ & 0.0997 & 0.0579 & $2.90 \pm 0.21$ \\
    534.41 & 5.81 & 5.52 & 5.52 & 5.44 & 6.30 & $5.72 \pm 0.16$ & 0.1049 & 0.0658 & $3.29 \pm 0.18$ \\
    \midrule
    \multicolumn{10}{c}{\textbf{Більший шків: $r_2 = \SI{30}{\milli\meter}$}} \\
    \midrule
    296.20 & 4.96 & 5.05 & 5.37 & 5.36 & 5.04 & $5.16 \pm 0.09$ & 0.0872 & 0.0809 & $2.70 \pm 0.09$ \\
    335.24 & 4.79 & 4.48 & 4.46 & 5.25 & 4.94 & $4.78 \pm 0.15$ & 0.0987 & 0.0941 & $3.14 \pm 0.20$ \\
    373.67 & 4.33 & 4.53 & 5.29 & 5.28 & 4.52 & $4.79 \pm 0.21$ & 0.1100 & 0.0938 & $3.13 \pm 0.27$ \\
    411.21 & 5.03 & 5.00 & 4.54 & 4.21 & 4.27 & $4.61 \pm 0.17$ & 0.1210 & 0.1013 & $3.38 \pm 0.26$ \\
    449.30 & 4.92 & 4.80 & 4.63 & 4.63 & 4.46 & $4.69 \pm 0.08$ & 0.1322 & 0.0979 & $3.26 \pm 0.11$ \\
    \bottomrule
  \end{tabular}
  \caption{Результати вимірювань та обчислень для Досліду 1 ($R_{\text{фікс}} = \SI{16.21}{\centi\meter}$, $h = \SI{1.076}{\meter}$)}
  \label{tab:exp1}
\end{table}

\subsection{Дослід 2. Залежність прискорення від розташування муфточок ($m = \text{const}$)}
При фіксованій масі падаючого вантажу $m_{\text{фікс}} \approx \SI{373}{\gram}$ (зокрема, $\SI{373.97}{\gram}$ для $r_1$ та $\SI{373.67}{\gram}$ для $r_2$, що суттєво перевищує порогову масу $m_{\text{повна}}$) змінювали положення чотирьох муфточок на спицях: $R \in \{5.0;\ 10.3;\ 16.0;\ 21.0;\ 23.5\}\ \si{\centi\meter}$. Для кожного положення було проведено по 5 спроб (по 2--3 паралельні відліки секундоміра в кожній). Результати зведено в Табл.~\ref{tab:exp2}.

\begin{table}[H]
  \centering
  \small
  \begin{tabular}{@{}cccccccccc@{}}
    \toprule
    $R$, \si{\centi\meter} & $\bar{t}_1$, \si{\second} & $\bar{t}_2$, \si{\second} & $\bar{t}_3$, \si{\second} & $\bar{t}_4$, \si{\second} & $\bar{t}_5$, \si{\second} & $\bar{t} \pm \Delta t$, \si{\second} & $(R/r)^2$ & $a$, \si{\meter\per\second\squared} & $g/a \pm \Delta(g/a)$ \\
    \midrule
    \multicolumn{10}{c}{\textbf{Менший шків: $r_1 = \SI{20}{\milli\meter}$ ($m = \SI{373.97}{\gram}$)}} \\
    \midrule
    5.0  & 5.66 & 6.32 & 6.33 & 5.43 & 5.78 & $5.90 \pm 0.18$ & 6.25   & 0.0618 & $159 \pm 10$ \\
    10.3 & 7.26 & 6.32 & 6.18 & 6.24 & 6.21 & $6.44 \pm 0.21$ & 26.52  & 0.0519 & $189 \pm 12$ \\
    16.0 & 7.05 & 8.66 & 10.29 & 9.53 & 10.64 & $9.23 \pm 0.64$ & 64.00  & 0.0253 & $388 \pm 54$ \\
    21.0 & 9.09 & 7.79 & 8.46 & 7.87 & 8.09 & $8.26 \pm 0.24$ & 110.25 & 0.0315 & $311 \pm 18$ \\
    23.5 & 9.98 & 8.46 & 7.96 & 7.90 & 7.69 & $8.40 \pm 0.42$ & 138.06 & 0.0305 & $321 \pm 32$ \\
    \midrule
    \multicolumn{10}{c}{\textbf{Більший шків: $r_2 = \SI{30}{\milli\meter}$ ($m = \SI{373.67}{\gram}$)}} \\
    \midrule
    5.0  & 3.76 & 3.55 & 3.55 & 3.56 & 3.75 & $3.63 \pm 0.05$ & 2.78   & 0.1631 & $60.1 \pm 1.7$ \\
    10.3 & 4.06 & 3.92 & 4.22 & 4.17 & 4.37 & $4.14 \pm 0.08$ & 11.79  & 0.1254 & $78.2 \pm 2.9$ \\
    16.0 & 4.33 & 4.53 & 5.29 & 5.28 & 4.52 & $4.79 \pm 0.21$ & 28.44  & 0.0938 & $104.5 \pm 9.0$ \\
    21.0 & 5.49 & 5.91 & 5.24 & 5.45 & 5.47 & $5.51 \pm 0.11$ & 49.00  & 0.0709 & $138.3 \pm 5.5$ \\
    23.5 & 5.61 & 5.44 & 5.49 & 5.66 & 5.57 & $5.55 \pm 0.04$ & 61.36  & 0.0699 & $140.4 \pm 2.1$ \\
    \bottomrule
  \end{tabular}
  \caption{Результати вимірювань та обчислень для Досліду 2 ($m_{\text{фікс}} \approx \SI{373}{\gram}$, $h = \SI{1.076}{\meter}$)}
  \label{tab:exp2}
\end{table}

% ==========================================================================
\section{Обробка результатів}
% ==========================================================================

\subsection{Оцінка похибок прямих і непрямих вимірювань}
Для кожної серії з $N = 5$ спроб обчислювали середнє арифметичне часу $\bar{t}$ та стандартну похибку середнього значення:
\begin{equation}
  \bar{t} = \frac{1}{N}\sum_{i=1}^{N} \bar{t}_i, \qquad \Delta t = \sqrt{\frac{\sum_{i=1}^{N}(\bar{t}_i - \bar{t})^2}{N(N-1)}}.
\end{equation}
Оскільки відносна похибка вимірювання висоти $\frac{\Delta h}{h} = \frac{0.001}{1.076} \approx 0.09\%$ та радіуса шківа є нехтовно малими порівняно з відносною похибкою часу $\frac{\Delta t}{\bar{t}} \sim 1\text{--}10\%$, похибки непрямих вимірювань прискорення $a$, кутового прискорення $\beta$ та величини $g/a$ визначаються формулами переносу похибок:
\begin{equation}
  \frac{\Delta a}{a} = \frac{\Delta \beta}{\beta} = \frac{\Delta(g/a)}{g/a} = \sqrt{\left(\frac{\Delta h}{h}\right)^2 + \left(2\frac{\Delta t}{\bar{t}}\right)^2} \approx 2\frac{\Delta t}{\bar{t}}.
\end{equation}

\subsection{Пряма оцінка моменту сили тертя}
За пороговими масами початку обертання маятника (Табл.~\ref{tab:constants}) маємо безпосередню оцінку моменту сили тертя $M_T \approx m_{\text{повна}} g r$:
\begin{itemize}
  \item Для шківа $r_1 = \SI{20}{\milli\meter}$: $M_{T,1} = \num{0.1408} \cdot \num{9.81} \cdot \num{0.020} = \SI{0.0276}{\newton\meter}$ (якщо ж враховувати лише додаткові важки $m_{0,1} = \SI{38.8}{\gram}$ при збалансованій платформі, то $M'_{T,1} = \SI{0.0076}{\newton\meter}$).
  \item Для шківа $r_2 = \SI{30}{\milli\meter}$: $M_{T,2} = \num{0.12824} \cdot \num{9.81} \cdot \num{0.030} = \SI{0.0377}{\newton\meter}$ (відповідно за додатковими важками $m_{0,2} = \SI{26.24}{\gram}$ маємо $M'_{T,2} = \SI{0.0077}{\newton\meter}$). Збіг значень $M'_{T,1} \approx M'_{T,2} \approx \SI{0.0077}{\newton\meter}$ свідчить про те, що власний момент сили тертя спокою в осі становить близько $\SI{7.7e-3}{\newton\meter}$.
\end{itemize}

\subsection{Аналіз Досліду 1: апроксимація залежності $\beta(mgr)$}
На Рис.~\ref{fig:exp1} зображено експериментальні точки та лінійні апроксимації $\beta = k \cdot (mgr) + b$, побудовані за методом найменших квадратів (МНК).

\begin{figure}[H]
  \centering
  \begin{tikzpicture}
    \begin{axis}[
      width=0.88\textwidth,
      height=7.2cm,
      xlabel={Момент зовнішньої сили $mgr$, \si{\newton\meter}},
      ylabel={Кутове прискорення $\beta$, \si{\radian\per\second\squared}},
      xmin=0.068, xmax=0.138,
      ymin=0.8, ymax=3.8,
      grid=both,
      grid style={line width=.1pt, draw=gray!20},
      major grid style={line width=.2pt, draw=gray!40},
      legend pos=south east,
      legend cell align={left},
      legend style={font=\small, fill=white, fill opacity=0.9}
    ]
      % Експериментальні точки r1 = 20 мм
      \addplot[
        only marks,
        color=series1,
        mark=*,
        mark size=2.2pt,
        error bars/.cd,
        y dir=both, y explicit
      ] coordinates {
        (0.0734, 1.263) +- (0, 0.176)
        (0.0807, 1.535) +- (0, 0.327)
        (0.0882, 2.337) +- (0, 0.169)
        (0.0997, 2.897) +- (0, 0.207)
        (0.1049, 3.289) +- (0, 0.182)
      };
      \addlegendentry{Експеримент ($r_1 = \SI{20}{\milli\meter}$)}

      % Лінія МНК r1 = 20 мм
      \addplot[
        domain=0.071:0.108,
        samples=50,
        color=series1,
        thick
      ] {65.725*x - 3.609};
      \addlegendentry{МНК ($r_1$): $\beta = 65.7(mgr) - 3.61$ ($R^2 = 0.982$)}

      % Експериментальні точки r2 = 30 мм
      \addplot[
        only marks,
        color=series2,
        mark=square*,
        mark size=2.2pt,
        error bars/.cd,
        y dir=both, y explicit
      ] coordinates {
        (0.0872, 2.698) +- (0, 0.090)
        (0.0987, 3.137) +- (0, 0.195)
        (0.1100, 3.128) +- (0, 0.268)
        (0.1210, 3.377) +- (0, 0.255)
        (0.1322, 3.262) +- (0, 0.108)
      };
      \addlegendentry{Експеримент ($r_2 = \SI{30}{\milli\meter}$)}

      % Лінія МНК r2 = 30 мм
      \addplot[
        domain=0.085:0.135,
        samples=50,
        color=series2,
        thick,
        dashed
      ] {12.209*x + 1.780};
      \addlegendentry{МНК ($r_2$): $\beta = 12.2(mgr) + 1.78$ ($R^2 = 0.711$)}
    \end{axis}
  \end{tikzpicture}
  \caption{Залежність кутового прискорення маятника $\beta$ від моменту сили тяжіння $mgr$ при $R_{\text{фікс}} = \SI{16.21}{\centi\meter}$}
  \label{fig:exp1}
\end{figure}

\begin{enumerate}
  \item \textbf{Менший шків ($r_1 = \SI{20}{\milli\meter}$):} Експериментальні точки з високою точністю лягають на пряму (коефіцієнт детермінації $R^2 = 0.9822$), що підтверджує справедливість рівняння моментів \eqref{eq:beta_linear}. Параметри МНК-апроксимації:
  \begin{equation*}
    k_1 = \frac{1}{I} = (\num{65.7} \pm \num{5.1})\ \si{\kilogram^{-1}\meter^{-2}}, \qquad b_1 = -\frac{M_T}{I} = (\num{-3.61} \pm \num{0.46})\ \si{\radian\per\second\squared}.
  \end{equation*}
  Звідси визначаємо повний момент інерції навантаженого маятника $I$, динамічний момент сили тертя $M_T$ (як точку перетину з віссю абсцис $M_T = -b_1/k_1$) та еквівалентну порогову масу $m_0 = \frac{M_T}{g r_1}$:
  \begin{equation}
    I = (\num{0.0152} \pm \num{0.0012})\ \si{\kilogram\meter\squared}, \qquad M_T = (\num{0.055} \pm \num{0.008})\ \si{\newton\meter}, \qquad m_0 = \SI{280}{\gram}.
  \end{equation}
  Той факт, що екстрапольоване значення $M_T$ перевищує момент тертя спокою ненавантаженої осі, пояснюється зростанням сили нормальної реакції в підшипниках під вагою великих вантажів ($m \sim 370\text{--}530\ \si{\gram}$) та втратами енергії на згинання і багатошарове розмотування нитки на малому шківі.

  \item \textbf{Більший шків ($r_2 = \SI{30}{\milli\meter}$):} Як видно з Рис.~\ref{fig:exp1}, починаючи з $m = \SI{335}{\gram}$ ($mgr \approx \SI{0.10}{\newton\meter}$), час падіння платформи виходить на плато ($\bar{t} \approx 4.6\text{--}4.8\ \si{\second}$), через що кутове прискорення майже перестає зростати ($\beta \approx 3.1\text{--}3.4\ \si{\radian\per\second\squared}$). Через це формальна лінійна регресія дає пологий нахил $k_2 = (\num{12.2} \pm \num{4.5})\ \si{\kilogram^{-1}\meter^{-2}}$ ($R^2 = 0.711$) та нефізичний додатний вільний член $b_2 = +1.78\ \si{\radian\per\second\squared}$. Причиною цього «насичення» при більшому плечі сили $r_2$ є проковзування витків нитки на шківі при більших прискореннях та суттєве зростання тертя в осі при перекосі шківа під дією великого моменту сили. Тому достовірним у Досліді 1 є результат для шківа $r_1$.
\end{enumerate}

\subsection{Аналіз Досліду 2: апроксимація залежності $g/a$ від $(R/r)^2$}
Згідно з рівнянням \eqref{eq:ga_linear}, графік залежності $y = g/a$ від $x = (R/r)^2$ має бути прямою лінією $y = k x + b$ з теоретичним кутовим коефіцієнтом:
\begin{equation}
  k_{\text{теор}} = \frac{2(m_s + m_b)}{m_{\text{фікс}}} = \frac{2(\num{100.77} + \num{125.34})}{\num{373.8}} = \frac{\num{452.22}}{\num{373.8}} \approx \num{1.21}.
  \label{eq:k_theor}
\end{equation}
Вільний член $b$ визначає момент інерції ненавантаженого маятника: $I_0 = b \cdot m_{\text{фікс}} r^2$. Графіки отриманих залежностей для обох шківів подано на Рис.~\ref{fig:exp2}.

\begin{figure}[H]
  \centering
  \begin{tikzpicture}
    \begin{axis}[
      width=0.88\textwidth,
      height=7.5cm,
      xlabel={Безрозмірний параметр $(R/r)^2$},
      ylabel={Відношення прискорень $g/a$},
      xmin=0, xmax=145,
      ymin=40, ymax=450,
      grid=both,
      grid style={line width=.1pt, draw=gray!20},
      major grid style={line width=.2pt, draw=gray!40},
      legend pos=north west,
      legend cell align={left},
      legend style={font=\small, fill=white, fill opacity=0.9}
    ]
      % Експериментальні точки r1 = 20 мм (основні 4 точки)
      \addplot[
        only marks,
        color=series1,
        mark=*,
        mark size=2.2pt,
        error bars/.cd,
        y dir=both, y explicit
      ] coordinates {
        (6.25, 158.8) +- (0, 9.7)
        (26.52, 189.1) +- (0, 12.0)
        (110.25, 310.9) +- (0, 17.9)
        (138.06, 321.4) +- (0, 31.8)
      };
      \addlegendentry{Експеримент ($r_1 = \SI{20}{\milli\meter}$)}

      % Викид r1 при R = 16 см
      \addplot[
        only marks,
        color=series1,
        mark=o,
        mark size=2.8pt,
        thick,
        error bars/.cd,
        y dir=both, y explicit
      ] coordinates {
        (64.00, 388.4) +- (0, 54.1)
      };
      \addlegendentry{Викид ($r_1$, $R = \SI{16}{\centi\meter}$)}

      % Лінія МНК r1 (без викиду)
      \addplot[
        domain=0:142,
        samples=50,
        color=series1,
        thick
      ] {1.293*x + 154.2};
      \addlegendentry{МНК ($r_1$): $g/a = 1.29(R/r)^2 + 154.2$ ($R^2 = 0.984$)}

      % Експериментальні точки r2 = 30 мм
      \addplot[
        only marks,
        color=series2,
        mark=square*,
        mark size=2.2pt,
        error bars/.cd,
        y dir=both, y explicit
      ] coordinates {
        (2.78, 60.1) +- (0, 1.7)
        (11.79, 78.2) +- (0, 2.9)
        (28.44, 104.5) +- (0, 9.0)
        (49.00, 138.3) +- (0, 5.5)
        (61.36, 140.4) +- (0, 2.1)
      };
      \addlegendentry{Експеримент ($r_2 = \SI{30}{\milli\meter}$)}

      % Лінія МНК r2 = 30 мм
      \addplot[
        domain=0:68,
        samples=50,
        color=series2,
        thick,
        dashed
      ] {1.430*x + 60.5};
      \addlegendentry{МНК ($r_2$): $g/a = 1.43(R/r)^2 + 60.5$ ($R^2 = 0.970$)}
    \end{axis}
  \end{tikzpicture}
  \caption{Залежність відношення $g/a$ від квадрата відносної відстані муфточок $(R/r)^2$ для двох шківів}
  \label{fig:exp2}
\end{figure}

\begin{enumerate}
  \item \textbf{Більший шків ($r_2 = \SI{30}{\milli\meter}$):} Усі 5 точок демонструють чітку лінійну залежність ($R^2 = 0.9700$):
  \begin{equation*}
    k_2 = \num{1.43} \pm \num{0.15}, \qquad b_2 = \num{60.5} \pm \num{5.5}.
  \end{equation*}
  Експериментальний нахил $k_2 = 1.43 \pm 0.15$ узгоджується з теоретичним значенням $k_{\text{теор}} = 1.21$. Момент інерції ненавантаженого маятника становить:
  \begin{equation}
    I_{0,2} = b_2 m_{\text{фікс}} r_2^2 = \num{60.5} \cdot \num{0.37367} \cdot (\num{0.030})^2 = (\num{0.0203} \pm \num{0.0018})\ \si{\kilogram\meter\squared}.
  \end{equation}

  \item \textbf{Менший шків ($r_1 = \SI{20}{\milli\meter}$):} Як видно з Табл.~\ref{tab:exp2}, у точці $R = \SI{16.0}{\centi\meter}$ під час 3--5 спроб сталося різке гальмування системи (час зріс із $\SI{7.05}{\second}$ у першій спробі до $\SI{10.64}{\second}$ у п'ятій, найімовірніше через тертя нитки об бортик шківа). Якщо враховувати цей викид, вільний член становить $b_1 = 188.8 \pm 60.5$ ($I_{0,1} = (\num{0.028} \pm \num{0.009})\ \si{\kilogram\meter\squared}$, $R^2 = 0.499$). Після виключення грубого промаху при $R = \SI{16.0}{\centi\meter}$ решта 4 точок лягають на пряму з високою точністю ($R^2 = 0.9835$):
  \begin{equation*}
    k_1 = \num{1.29} \pm \num{0.12}, \qquad b_1 = \num{154.2} \pm \num{10.6}.
  \end{equation*}
  Отриманий кутовий коефіцієнт $k_1 = 1.29 \pm 0.12$ у межах похибки збігається з теоретичним відношенням мас $k_{\text{теор}} = 1.21$, а власний момент інерції маятника дорівнює:
  \begin{equation}
    I_{0,1} = b_1 m_{\text{фікс}} r_1^2 = \num{154.2} \cdot \num{0.37397} \cdot (\num{0.020})^2 = (\num{0.0231} \pm \num{0.0016})\ \si{\kilogram\meter\squared}.
  \end{equation}
\end{enumerate}

Усереднюючи значення $I_{0,1}$ та $I_{0,2}$, одержуємо підсумкове значення моменту інерції ненавантаженого маятника Обербека:
\begin{equation}
  I_0 = (\num{0.0217} \pm \num{0.0017})\ \si{\kilogram\meter\squared}.
\end{equation}

\begin{keyfacts}
\textbf{Основні результати:} \\
Сумарний момент інерції навантаженого маятника (Дослід 1, $R = \SI{16.21}{\centi\meter}$): $I = (\num{0.0152} \pm \num{0.0012})\ \si{\kilogram\meter\squared}$ (відносна похибка $\varepsilon = 7.9\%$). \\
Момент інерції ненавантаженого маятника (Дослід 2): $I_0 = (\num{0.0217} \pm \num{0.0017})\ \si{\kilogram\meter\squared}$ (відносна похибка $\varepsilon = 7.8\%$). \\
Експериментальний нахил $g/a$ від $(R/r)^2$: $k_1 = \num{1.29} \pm \num{0.12}$ та $k_2 = \num{1.43} \pm \num{0.15}$ при теоретичному значенні $k_{\text{теор}} = \num{1.21}$.
\end{keyfacts}

% ==========================================================================
\section{Висновки}
% ==========================================================================
\begin{enumerate}
  \item У ході лабораторної роботи досліджено динаміку обертального руху твердого тіла на прикладі маятника Обербека та двома незалежними методами підтверджено лінійний характер рівняння моментів відносно нерухомої осі.
  \item У першому досліді при сталому моменті інерції ($R_{\text{фікс}} = \SI{16.21}{\centi\meter}$) для меншого шківа ($r_1 = \SI{20}{\milli\meter}$) експериментальна залежність кутового прискорення $\beta$ від моменту зовнішньої сили $mgr$ є строго лінійною ($R^2 = 0.982$), звідки знайдено момент інерції навантаженого маятника $I = (\num{0.0152} \pm \num{0.0012})\ \si{\kilogram\meter\squared}$.
  \item У другому досліді при зміні відстані муфточок $R$ від осі обертання підтверджено квадратичну залежність моменту інерції від відстані муфт (теорему Штейнера): графіки залежності $g/a$ від $(R/r)^2$ є лінійними ($R^2 = 0.984$ для $r_1$ та $R^2 = 0.970$ for $r_2$). Експериментальні кутові коефіцієнти прямих ($k_1 = 1.29 \pm 0.12$ та $k_2 = 1.43 \pm 0.15$) у межах похибки узгоджуються з теоретичним відношенням сумарної маси чотирьох муфточок до маси падаючого вантажу $k_{\text{теор}} = 1.21$. Значення власного моменту інерції ненавантаженого маятника, отримані на двох різних шківах ($I_{0,1} = \SI{0.0231}{\kilogram\meter\squared}$ та $I_{0,2} = \SI{0.0203}{\kilogram\meter\squared}$), узгоджуються між собою в межах стандартної похибки і дають середнє значення $I_0 = (\num{0.0217} \pm \num{0.0017})\ \si{\kilogram\meter\squared}$.
  \item Деяка розбіжність між значенням повного моменту інерції $I$ з першого досліду та величиною $I_0 + 2(m_s+m_b)R^2 \approx \SI{0.0336}{\kilogram\meter\squared}$ із другого досліду, а також ефект «насичення» прискорення на більшому шківі $r_2$ у першому досліді пояснюються систематичними факторами:
  \begin{itemize}
    \item нехтуванням моментом сили тертя $M_T$ при виведенні формули \eqref{eq:ga_linear} для другого досліду (оскільки наявність тертя зменшує фактичне прискорення $a$, відношення $g/a$ систематично завищується, що призводить до завищеної оцінки вільного члена $b$ і моменту інерції $I_0$);
    \item зміною ефективного радіуса намотування нитки (накладанням кількох шарів витків на вузькому шківі) та тертям нитки об бічні стінки шківа під час руху (що особливо проявилося у вигляді статистичного викиду при $R = \SI{16}{\centi\meter}$ на шківі $r_1$);
    \item залежністю моменту сили тертя в підшипниках від маси підвішеного вантажу та людським фактором (часом реакції при ручному запуску секундомірів).
  \end{itemize}
\end{enumerate}

\end{document}